\section{Discussion}
\label{sec:discussion}

\paragraph{What the known-answer study shows}
The world model's explanation is faithful by construction: the baseline does not
move with the plan, the lift is zero when nothing is planned, and the two add up
to the forecast. In the known-answer study that explanation was right about what
each lever adds when plans were made at random, and wrong in two ways that no
check against the model could see. When merchants plan into strong weeks, the
lifts are \kaCoWmBiasLo{} to \kaCoWmBiasHi{} times the truth, because the data
tie the plan to demand the model cannot see. And asked about a change in the
regular price, the model states a response near the temporary-cut elasticity,
\kaClWmTimesLo{} to \kaCoWmTimesHi{} times the truth. Molecular attributions in
MolSanity~\cite{molsanity} came apart from the truth for a similar reason: a
model can be faithful to what it learned and what it learned can be wrong.

\paragraph{What this means for a merchant reading a lift}
In the known-answer study, a lift read from a model fitted to the store's own
history inherited its planning habits. If promotions go into weeks that would
sell well anyway, the model's promotion lift overstates what the promotion adds,
and an accurate forecast is no protection: the correct world-model seeds and
LightGBM forecast better when plans followed demand than when they were made at
random. The check against what sold, the one most often run in practice, agrees
with the overstated lift for the same reason. Two things follow. Agreement
between models fitted to the same data is not confirmation. And a model's
answer to a question it was not trained on, here a lasting change in the
regular price from a model that mostly saw temporary cuts, needs checking
against a measured response before it is used to set prices.

\paragraph{Which checks to run}
The known-answer study ranks the checks. The price-direction check (C4) is cheap
and caught every model blind to price, though it also flags LightGBM, whose
response to price comes in steps. Per-lever agreement with a second model, an
exploratory comparison, caught lifts given to the wrong lever, which every
check on the total lift missed. Construction (C1) and seed stability (C3) are
worth confirming once. The check against what sold (C5) and agreement with
LightGBM on the total lift (C6) confirm agreement, not correctness.
Randomisation (C2), the standard sanity check for
attributions~\cite{adebayo2018}, does not test an explanation built into the
model: re-initialising every part left the lifts' ranking largely in place.
Versions of C2 on the lift as a share of the baseline, or within planned weeks
only, were not tested. A measured cause (C7) is the only check that found the
regular-price error, though with six spells it could not tell a missing price
response from a correct one. No check in this study finds the bias from
planning into strong weeks; a measured promotion effect, which public retail
data do not offer, was not tested.

\paragraph{Measured causes on public data are scarce and noisy}
Dominick's zone price deviations give a pooled 8-week response close to a
placebo at false dates and a full-episode response in line with the chain's
own experiments; intervals clustered by a handful of stores are far too narrow,
and prices start moving before the detector sees them. A retailer running its
own tests would do better with the design under its control: stores assigned at
random, a known start date, and enough stores per arm for honest intervals.

\paragraph{The audits of the companion paper's models}
On the final tuned world models, the explanations add up under each model's own
plan for nothing planned on every window, the copies agree, and the state
carries the plan. Two results matter to a merchant. First, on Dominick's
and M5 a deeper cut can lower the forecast, a fault no accuracy score
revealed and one a merchant setting prices from the model would meet directly;
the earlier audit found the same on Dominick's and M5. Second, both the world
model and LightGBM state about twice the measured response to the zone price
deviations. In the known-answer study only the world model overstated this
response, so the real data do not reproduce that pattern, and the measured
interval is too wide to say which model is closer. The earlier audit also showed
that C1 under the audit's plan measures the distance between two definitions of
nothing planned, and that for an item with no recent price the model computes
its baseline at its mean price or at 1.0, which no merchant would call nothing
planned. An audit has to read the model's own definition of nothing planned,
and check it, before C1 means anything. Across the searches' versions, passing
more checks was not related to accuracy. For a retailer this means two separate
tests: an accuracy comparison to choose the forecast, and checks of the
explanation, with a price-direction test and a measured cause among them,
before its lifts are used to plan.

\section{Conclusion}
\label{sec:conclusion}

Plan-lift explanations built into a forecasting model can be faithful, stable
and in agreement with sales and with another model, and still be wrong about
what a plan causes. In generated store data with known effects, the checks that
compare an explanation with the model, the data or another model caught a model
blind to price and, compared lever by lever, lifts given to the wrong lever.
They did not catch lifts inflated by planning into strong weeks, or a model
whose response to a regular price change is near its response to a temporary
cut. Only a measured cause found the second. On public retail data the one such
cause we could measure, Dominick's zone price deviations, gives a pooled
response near the chain's experiments over the full episode, but its 8-week
estimate is close to a placebo at false dates, so on public data the comparison
is descriptive. The final tuned world models of the companion paper add up
under their own plan for nothing planned on every window, but on most data sets
a deeper cut can lower their forecast, and both they and LightGBM state about
twice the measured response to a regular price change. Across the
configurations the settings searches tried, the number of checks passed was not
related to how well they forecast. On these models and data sets, an accurate
forecast was not evidence of a correct explanation, and the explanation has to
be tested on its own.
